Welcome to the proceedings of the Shared Tasks at the \textbf{Fourth Arabic Natural Language Processing Conference (ArabicNLP 2026)}, co-located with \textbf{EMNLP 2026} in Budapest, Hungary, October 28 and 29, 2026.

This volume represents a major achievement in collaborative Arabic NLP research, bringing together a total of \textbf{212 papers from 15 shared tasks}, including \textbf{15 overview papers} and \textbf{197 system description papers} from participating teams worldwide. Each of the 15 overview papers and each system description paper was reviewed by three reviewers. This substantial collection reflects the growing vitality and maturity of the Arabic NLP research community and showcases the field's expansion into increasingly diverse and sophisticated challenges.

This year marks a new record for ArabicNLP, featuring the largest number of shared tasks to date, \textbf{15 in total}, up from 11 last year, and the largest number of system description papers, up from 127. The collection highlights the community's growing interest in evaluating, aligning and extending large language models for Arabic across text, speech and multimodal domains, as well as in culturally grounded and ethically sensitive applications. The process began with \textbf{27 submitted shared task proposals}, of which \textbf{20 were accepted}: \textbf{10 proceeded independently}, while \textbf{10 were merged into five tasks}, \textbf{resulting in 15 shared tasks in total}. This outcome underscores the community's strong spirit of collaboration and its collective effort to design impactful, high quality shared tasks.

\subsection*{\textbf{Organization and Participation of the Shared Tasks}}

The 15 shared tasks at ArabicNLP 2026 are organized into four thematic tracks, each addressing critical needs and emerging priorities in Arabic language technology. Together, they represent the community's most comprehensive shared task effort to date, covering speech and multimodal processing, dialectal variation, text structure and readability, content safety, and the evaluation of knowledge, hallucination and agentic behavior in large language models.

The paper counts below refer to system description papers; each task also contributes one overview paper.

\subsubsection*{\textbf{Track 1: Speech, Multimodal and Dialectal Processing}}

This track features five shared tasks that advance Arabic language processing beyond standard written text:

\begin{itemize}
\item \textbf{NADI 2026} (Robust and Mixed-Dialect Arabic Speech Processing), 14 papers  
\item \textbf{ArA-DF 2026} (Arabic Speech Deepfake Detection), 8 papers  
\item \textbf{ImageEval 2026} (Culturally Grounded Arabic Multimodal Evaluation), 12 papers  
\item \textbf{AlexandriaX 2026} (Context-Aware Dialectal Arabic Machine Translation and Evaluation), 12 papers  
\item \textbf{DialectSentEval 2026} (Arabic Dialect Sentiment Analysis and Dialect Swapping), 16 papers
\end{itemize}

These tasks address the pressing need for technologies that can process Arabic across different modalities and the full range of dialectal variation.

\subsubsection*{\textbf{Track 2: Text Structure, Readability and Discourse}}

This track comprises four shared tasks focused on the structure and accessibility of Arabic text:

\begin{itemize}
\item \textbf{BAREC 2026} (Second Shared Task on Arabic Readability Assessment), 10 papers  
\item \textbf{Arabic Sentence Segmentation}, 14 papers  
\item \textbf{Daleel 2026} (Arabic Argumentative Discourse Mining), 10 papers  
\item \textbf{AraGenre} (Hierarchical Definition-Guided Arabic Genre Classification), 11 papers
\end{itemize}

These tasks tackle foundational challenges in automated assessment, text segmentation and discourse analysis, where Arabic has long lacked standardized benchmarks.

\subsubsection*{\textbf{Track 3: Safety, Stance and Content Moderation}}

This track presents two shared tasks focused on harmful content detection and stance classification in Arabic.

\begin{itemize}
\item \textbf{StanceEval 2026} (Stance Detection in Arabic), 18 papers  
\item \textbf{ArGuard} (Detection of Harmful Content in Arabic Memes and LLM Prompts), 25 papers
\end{itemize}

These tasks respond to the urgent practical need for moderation and safety technologies that work reliably on Arabic content, including multimodal content and adversarial prompts.

\subsubsection*{\textbf{Track 4: Knowledge, Hallucination and Agentic Evaluation}}

This track introduces four shared tasks designed to evaluate what large language models know, when they fabricate, and how they act:

\begin{itemize}
\item \textbf{IslamicEval 2026} (Fine-Grained Hallucination Detection in Arabic Islamic Content), 11 papers  
\item \textbf{HalluScoring} (LLM Hallucination Detection in Arabic Question Answering), 10 papers  
\item \textbf{KnowledgeGraphEval 2026} (Arabic Knowledge Graph Construction and Domain Adaptation), 12 papers  
\item \textbf{AISA-ArabicFC} (Arabic Function Calling for Agentic AI Systems), 14 papers
\end{itemize}

These tasks represent a critical step toward ensuring that AI systems represent Arabic knowledge accurately, signal their own uncertainty, and behave reliably when given the ability to act.

\subsection*{\textbf{ACommunity Impact and Participation}}

The response to this year's shared tasks, with \textbf{197 system description papers across the 15 tasks}, demonstrates the Arabic NLP community's continued growth and dynamism. Participating teams represent a diverse range of institutions across multiple continents, including universities, research centers and industry partners, reflecting both the global interest in Arabic NLP and the real-world relevance of these challenges.

The breadth of methodological approaches presented in these proceedings is particularly noteworthy, spanning traditional machine learning techniques, state-of-the-art transformer models, retrieval-augmented generation systems, multimodal architectures and agentic pipelines. This methodological diversity reflects the field's technical maturity and the creativity researchers bring to the specific challenges of Arabic language processing.

Select system description papers from this volume are presented in the two poster sessions of the conference, and each shared task is introduced to the full audience in a highlights segment at the end of an oral session, seven tasks on the first day and eight on the second.

\subsection*{\textbf{Acknowledgments}}

We extend our sincere gratitude to the dedicated organizers of each shared task. Their expertise in task design, dataset development, evaluation framework creation and ongoing participant support has been fundamental to the success of these shared tasks. The quality and impact of this volume are a direct reflection of their commitment and hard work.

We thank all participating teams whose innovative research pushes the boundaries of Arabic NLP. Their contributions address not only technical challenges but also important societal needs in education, cultural preservation, content moderation and ethical AI development.

We also acknowledge the program committee members, reviewers, our fellow chairs and the ArabicNLP 2026 organizing committee for their invaluable support throughout the shared task cycle.

Finally, we gratefully acknowledge our sponsors, whose support has made this conference and these proceedings possible.

We hope this volume serves as both a comprehensive record of the current state of Arabic NLP research and a foundation for future innovations that will continue to advance language technologies for Arabic-speaking communities worldwide.

\textbf{Sakhar Alkhereyf}, HUMAIN, Saudi Arabia.

\textbf{Hamzah Luqman}, King Fahd University of Petroleum and Minerals, Saudi Arabia.

\textbf{Shared Tasks Chairs}


Website of the shared Tasks: \href{https://arabicnlp2026.sigarab.org/shared-tasks}{arabicnlp2026.sigarab.org/shared-tasks}
